% Options for packages loaded elsewhere
\PassOptionsToPackage{unicode}{hyperref}
\PassOptionsToPackage{hyphens}{url}
\documentclass[
]{article}
\usepackage{xcolor}
\usepackage[margin=1in]{geometry}
\usepackage{amsmath,amssymb}
\setcounter{secnumdepth}{-\maxdimen} % remove section numbering
\usepackage{iftex}
\ifPDFTeX
  \usepackage[T1]{fontenc}
  \usepackage[utf8]{inputenc}
  \usepackage{textcomp} % provide euro and other symbols
\else % if luatex or xetex
  \usepackage{unicode-math} % this also loads fontspec
  \defaultfontfeatures{Scale=MatchLowercase}
  \defaultfontfeatures[\rmfamily]{Ligatures=TeX,Scale=1}
\fi
\usepackage{lmodern}
\ifPDFTeX\else
  % xetex/luatex font selection
\fi
% Use upquote if available, for straight quotes in verbatim environments
\IfFileExists{upquote.sty}{\usepackage{upquote}}{}
\IfFileExists{microtype.sty}{% use microtype if available
  \usepackage[]{microtype}
  \UseMicrotypeSet[protrusion]{basicmath} % disable protrusion for tt fonts
}{}
\makeatletter
\@ifundefined{KOMAClassName}{% if non-KOMA class
  \IfFileExists{parskip.sty}{%
    \usepackage{parskip}
  }{% else
    \setlength{\parindent}{0pt}
    \setlength{\parskip}{6pt plus 2pt minus 1pt}}
}{% if KOMA class
  \KOMAoptions{parskip=half}}
\makeatother
\usepackage{color}
\usepackage{fancyvrb}
\newcommand{\VerbBar}{|}
\newcommand{\VERB}{\Verb[commandchars=\\\{\}]}
\DefineVerbatimEnvironment{Highlighting}{Verbatim}{commandchars=\\\{\}}
% Add ',fontsize=\small' for more characters per line
\usepackage{framed}
\definecolor{shadecolor}{RGB}{248,248,248}
\newenvironment{Shaded}{\begin{snugshade}}{\end{snugshade}}
\newcommand{\AlertTok}[1]{\textcolor[rgb]{0.94,0.16,0.16}{#1}}
\newcommand{\AnnotationTok}[1]{\textcolor[rgb]{0.56,0.35,0.01}{\textbf{\textit{#1}}}}
\newcommand{\AttributeTok}[1]{\textcolor[rgb]{0.13,0.29,0.53}{#1}}
\newcommand{\BaseNTok}[1]{\textcolor[rgb]{0.00,0.00,0.81}{#1}}
\newcommand{\BuiltInTok}[1]{#1}
\newcommand{\CharTok}[1]{\textcolor[rgb]{0.31,0.60,0.02}{#1}}
\newcommand{\CommentTok}[1]{\textcolor[rgb]{0.56,0.35,0.01}{\textit{#1}}}
\newcommand{\CommentVarTok}[1]{\textcolor[rgb]{0.56,0.35,0.01}{\textbf{\textit{#1}}}}
\newcommand{\ConstantTok}[1]{\textcolor[rgb]{0.56,0.35,0.01}{#1}}
\newcommand{\ControlFlowTok}[1]{\textcolor[rgb]{0.13,0.29,0.53}{\textbf{#1}}}
\newcommand{\DataTypeTok}[1]{\textcolor[rgb]{0.13,0.29,0.53}{#1}}
\newcommand{\DecValTok}[1]{\textcolor[rgb]{0.00,0.00,0.81}{#1}}
\newcommand{\DocumentationTok}[1]{\textcolor[rgb]{0.56,0.35,0.01}{\textbf{\textit{#1}}}}
\newcommand{\ErrorTok}[1]{\textcolor[rgb]{0.64,0.00,0.00}{\textbf{#1}}}
\newcommand{\ExtensionTok}[1]{#1}
\newcommand{\FloatTok}[1]{\textcolor[rgb]{0.00,0.00,0.81}{#1}}
\newcommand{\FunctionTok}[1]{\textcolor[rgb]{0.13,0.29,0.53}{\textbf{#1}}}
\newcommand{\ImportTok}[1]{#1}
\newcommand{\InformationTok}[1]{\textcolor[rgb]{0.56,0.35,0.01}{\textbf{\textit{#1}}}}
\newcommand{\KeywordTok}[1]{\textcolor[rgb]{0.13,0.29,0.53}{\textbf{#1}}}
\newcommand{\NormalTok}[1]{#1}
\newcommand{\OperatorTok}[1]{\textcolor[rgb]{0.81,0.36,0.00}{\textbf{#1}}}
\newcommand{\OtherTok}[1]{\textcolor[rgb]{0.56,0.35,0.01}{#1}}
\newcommand{\PreprocessorTok}[1]{\textcolor[rgb]{0.56,0.35,0.01}{\textit{#1}}}
\newcommand{\RegionMarkerTok}[1]{#1}
\newcommand{\SpecialCharTok}[1]{\textcolor[rgb]{0.81,0.36,0.00}{\textbf{#1}}}
\newcommand{\SpecialStringTok}[1]{\textcolor[rgb]{0.31,0.60,0.02}{#1}}
\newcommand{\StringTok}[1]{\textcolor[rgb]{0.31,0.60,0.02}{#1}}
\newcommand{\VariableTok}[1]{\textcolor[rgb]{0.00,0.00,0.00}{#1}}
\newcommand{\VerbatimStringTok}[1]{\textcolor[rgb]{0.31,0.60,0.02}{#1}}
\newcommand{\WarningTok}[1]{\textcolor[rgb]{0.56,0.35,0.01}{\textbf{\textit{#1}}}}
\usepackage{graphicx}
\makeatletter
\newsavebox\pandoc@box
\newcommand*\pandocbounded[1]{% scales image to fit in text height/width
  \sbox\pandoc@box{#1}%
  \Gscale@div\@tempa{\textheight}{\dimexpr\ht\pandoc@box+\dp\pandoc@box\relax}%
  \Gscale@div\@tempb{\linewidth}{\wd\pandoc@box}%
  \ifdim\@tempb\p@<\@tempa\p@\let\@tempa\@tempb\fi% select the smaller of both
  \ifdim\@tempa\p@<\p@\scalebox{\@tempa}{\usebox\pandoc@box}%
  \else\usebox{\pandoc@box}%
  \fi%
}
% Set default figure placement to htbp
\def\fps@figure{htbp}
\makeatother
\setlength{\emergencystretch}{3em} % prevent overfull lines
\providecommand{\tightlist}{%
  \setlength{\itemsep}{0pt}\setlength{\parskip}{0pt}}
\usepackage{bookmark}
\IfFileExists{xurl.sty}{\usepackage{xurl}}{} % add URL line breaks if available
\urlstyle{same}
\hypersetup{
  pdftitle={R Programming Practical Examination},
  pdfauthor={Jerjin Joy Doronila},
  hidelinks,
  pdfcreator={LaTeX via pandoc}}

\title{R Programming Practical Examination}
\author{Jerjin Joy Doronila}
\date{2026-11-11}

\begin{document}
\maketitle

\subsubsection{Part A: Basic Operations}\label{part-a-basic-operations}

\textbf{A1} Create v as the sequence 5, 10, 15,\ldots, 50 using
\texttt{seq()}. Print its length, sum and mean.

\begin{Shaded}
\begin{Highlighting}[]
\NormalTok{v }\OtherTok{\textless{}{-}} \FunctionTok{seq}\NormalTok{(}\DecValTok{5}\NormalTok{, }\DecValTok{50}\NormalTok{, }\AttributeTok{by =} \DecValTok{5}\NormalTok{)}
\NormalTok{v}
\end{Highlighting}
\end{Shaded}

\begin{verbatim}
##  [1]  5 10 15 20 25 30 35 40 45 50
\end{verbatim}

\begin{Shaded}
\begin{Highlighting}[]
\FunctionTok{length}\NormalTok{(v)}
\end{Highlighting}
\end{Shaded}

\begin{verbatim}
## [1] 10
\end{verbatim}

\begin{Shaded}
\begin{Highlighting}[]
\FunctionTok{sum}\NormalTok{(v)}
\end{Highlighting}
\end{Shaded}

\begin{verbatim}
## [1] 275
\end{verbatim}

\begin{Shaded}
\begin{Highlighting}[]
\FunctionTok{mean}\NormalTok{(v)}
\end{Highlighting}
\end{Shaded}

\begin{verbatim}
## [1] 27.5
\end{verbatim}

\textbf{A2} Create w that repeats the vector \texttt{c(1,\ 2,\ 3)} four
times, and z that repeats each element of \texttt{c(1,\ 2,\ 3)} four
times. Print both.

\begin{Shaded}
\begin{Highlighting}[]
\NormalTok{w }\OtherTok{\textless{}{-}} \FunctionTok{c}\NormalTok{(}\DecValTok{1}\NormalTok{, }\DecValTok{2}\NormalTok{, }\DecValTok{3}\NormalTok{)}
\NormalTok{w }\OtherTok{\textless{}{-}} \FunctionTok{rep}\NormalTok{(w, }\DecValTok{4}\NormalTok{)}

\NormalTok{z }\OtherTok{\textless{}{-}} \FunctionTok{c}\NormalTok{(}\DecValTok{1}\NormalTok{, }\DecValTok{2}\NormalTok{, }\DecValTok{3}\NormalTok{)}
\NormalTok{z }\OtherTok{\textless{}{-}} \FunctionTok{rep}\NormalTok{(z, }\AttributeTok{each =} \DecValTok{4}\NormalTok{)}

\NormalTok{w}
\end{Highlighting}
\end{Shaded}

\begin{verbatim}
##  [1] 1 2 3 1 2 3 1 2 3 1 2 3
\end{verbatim}

\begin{Shaded}
\begin{Highlighting}[]
\NormalTok{z}
\end{Highlighting}
\end{Shaded}

\begin{verbatim}
##  [1] 1 1 1 1 2 2 2 2 3 3 3 3
\end{verbatim}

\textbf{A3} Compute and print: the remainder of 17 divided by 5, the
integer division of 17 by 5, 2 to the power 10, the square root of 144,
and pi rounded to 3 decimals.

\begin{Shaded}
\begin{Highlighting}[]
\DecValTok{17} \SpecialCharTok{\%\%} \DecValTok{5}
\end{Highlighting}
\end{Shaded}

\begin{verbatim}
## [1] 2
\end{verbatim}

\begin{Shaded}
\begin{Highlighting}[]
\DecValTok{17} \SpecialCharTok{\%/\%} \DecValTok{5}
\end{Highlighting}
\end{Shaded}

\begin{verbatim}
## [1] 3
\end{verbatim}

\begin{Shaded}
\begin{Highlighting}[]
\DecValTok{2}\SpecialCharTok{\^{}}\DecValTok{10}
\end{Highlighting}
\end{Shaded}

\begin{verbatim}
## [1] 1024
\end{verbatim}

\begin{Shaded}
\begin{Highlighting}[]
\FunctionTok{sqrt}\NormalTok{(}\DecValTok{144}\NormalTok{)}
\end{Highlighting}
\end{Shaded}

\begin{verbatim}
## [1] 12
\end{verbatim}

\begin{Shaded}
\begin{Highlighting}[]
\NormalTok{pi }\OtherTok{\textless{}{-}} \FloatTok{3.1415926535897932384626433832795}
\FunctionTok{round}\NormalTok{(pi, }\AttributeTok{digits =} \DecValTok{3}\NormalTok{)}
\end{Highlighting}
\end{Shaded}

\begin{verbatim}
## [1] 3.142
\end{verbatim}

\textbf{A4} Given a \textless- c(2, 4, 6, 8) and b \textless- c(1, 3, 5,
7), print a + b, a * b, the logical comparison a \textgreater{} b, and
the dot product sum(a*b).

\begin{Shaded}
\begin{Highlighting}[]
\NormalTok{a }\OtherTok{\textless{}{-}} \FunctionTok{c}\NormalTok{(}\DecValTok{2}\NormalTok{, }\DecValTok{4}\NormalTok{, }\DecValTok{6}\NormalTok{, }\DecValTok{8}\NormalTok{)}
\NormalTok{b }\OtherTok{\textless{}{-}} \FunctionTok{c}\NormalTok{(}\DecValTok{1}\NormalTok{, }\DecValTok{3}\NormalTok{, }\DecValTok{5}\NormalTok{, }\DecValTok{7}\NormalTok{)}

\NormalTok{a }\SpecialCharTok{+}\NormalTok{ b}
\end{Highlighting}
\end{Shaded}

\begin{verbatim}
## [1]  3  7 11 15
\end{verbatim}

\begin{Shaded}
\begin{Highlighting}[]
\NormalTok{a }\SpecialCharTok{*}\NormalTok{ b}
\end{Highlighting}
\end{Shaded}

\begin{verbatim}
## [1]  2 12 30 56
\end{verbatim}

\begin{Shaded}
\begin{Highlighting}[]
\NormalTok{a }\SpecialCharTok{\textgreater{}}\NormalTok{ b}
\end{Highlighting}
\end{Shaded}

\begin{verbatim}
## [1] TRUE TRUE TRUE TRUE
\end{verbatim}

\begin{Shaded}
\begin{Highlighting}[]
\FunctionTok{sum}\NormalTok{(a}\SpecialCharTok{*}\NormalTok{b)}
\end{Highlighting}
\end{Shaded}

\begin{verbatim}
## [1] 100
\end{verbatim}

\textbf{A5} Use paste() to produce ``Student-1''
``Student-2''Student-3''. Then identify the datatype of 42, 42L, ``42'',
41, and TRUE. what are the inputs.

\begin{Shaded}
\begin{Highlighting}[]
\FunctionTok{paste}\NormalTok{(}\StringTok{"Student{-}1"}\NormalTok{,}\StringTok{"Student{-}2"}\NormalTok{, }\StringTok{"Student{-}3"}\NormalTok{)}
\end{Highlighting}
\end{Shaded}

\begin{verbatim}
## [1] "Student-1 Student-2 Student-3"
\end{verbatim}

\begin{Shaded}
\begin{Highlighting}[]
\FunctionTok{typeof}\NormalTok{(}\DecValTok{42}\NormalTok{)}
\end{Highlighting}
\end{Shaded}

\begin{verbatim}
## [1] "double"
\end{verbatim}

\begin{Shaded}
\begin{Highlighting}[]
\FunctionTok{typeof}\NormalTok{(}\DecValTok{42}\DataTypeTok{L}\NormalTok{)}
\end{Highlighting}
\end{Shaded}

\begin{verbatim}
## [1] "integer"
\end{verbatim}

\begin{Shaded}
\begin{Highlighting}[]
\FunctionTok{typeof}\NormalTok{(}\StringTok{"42"}\NormalTok{)}
\end{Highlighting}
\end{Shaded}

\begin{verbatim}
## [1] "character"
\end{verbatim}

\begin{Shaded}
\begin{Highlighting}[]
\FunctionTok{typeof}\NormalTok{(}\DecValTok{41}\DataTypeTok{i}\NormalTok{)}
\end{Highlighting}
\end{Shaded}

\begin{verbatim}
## [1] "complex"
\end{verbatim}

\begin{Shaded}
\begin{Highlighting}[]
\FunctionTok{typeof}\NormalTok{(}\ConstantTok{TRUE}\NormalTok{)}
\end{Highlighting}
\end{Shaded}

\begin{verbatim}
## [1] "logical"
\end{verbatim}

\subsubsection{Part B: Data Frames and Conditonal
Statements}\label{part-b-data-frames-and-conditonal-statements}

\textbf{B1} Create a data frame \texttt{students} with 8 rows and these
columns (use \texttt{stringAsFactor\ =\ FALSE}, sex is ``M''or ``F''):

\begin{Shaded}
\begin{Highlighting}[]
\NormalTok{students }\OtherTok{\textless{}{-}} \FunctionTok{data.frame}\NormalTok{(}
  \AttributeTok{id =} \FunctionTok{c}\NormalTok{(}\DecValTok{1}\SpecialCharTok{:}\DecValTok{8}\NormalTok{),}
  \AttributeTok{name =} \FunctionTok{c}\NormalTok{(}\StringTok{"Alice"}\NormalTok{, }\StringTok{"Bruno"}\NormalTok{, }\StringTok{"Chen"}\NormalTok{, }\StringTok{"Dana"}\NormalTok{, }\StringTok{"Elif"}\NormalTok{, }\StringTok{"Farid"}\NormalTok{, }\StringTok{"Grace"}\NormalTok{, }\StringTok{"Hiro"}\NormalTok{),}
  \AttributeTok{sex =} \FunctionTok{c}\NormalTok{(}\StringTok{"F"}\NormalTok{, }\StringTok{"M"}\NormalTok{,}\StringTok{"M"}\NormalTok{,  }\StringTok{"F"}\NormalTok{, }\StringTok{"F"}\NormalTok{, }\StringTok{"M"}\NormalTok{, }\StringTok{"F"}\NormalTok{, }\StringTok{"M"}\NormalTok{), }\AttributeTok{stringAsFactor =} \ConstantTok{FALSE}\NormalTok{, }
  \AttributeTok{section =} \FunctionTok{c}\NormalTok{(}\StringTok{"A"}\NormalTok{, }\StringTok{"B"}\NormalTok{, }\StringTok{"A"}\NormalTok{, }\StringTok{"C"}\NormalTok{, }\StringTok{"B"}\NormalTok{, }\StringTok{"C"}\NormalTok{, }\StringTok{"A"}\NormalTok{, }\StringTok{"B"}\NormalTok{),}
  \AttributeTok{midterm =} \FunctionTok{c}\NormalTok{(}\DecValTok{78}\NormalTok{, }\DecValTok{85}\NormalTok{, }\DecValTok{92}\NormalTok{, }\DecValTok{67}\NormalTok{, }\DecValTok{74}\NormalTok{, }\DecValTok{88}\NormalTok{, }\DecValTok{95}\NormalTok{, }\DecValTok{59}\NormalTok{),}
  \AttributeTok{final =} \FunctionTok{c}\NormalTok{(}\DecValTok{82}\NormalTok{, }\DecValTok{79}\NormalTok{, }\DecValTok{96}\NormalTok{, }\DecValTok{71}\NormalTok{, }\DecValTok{80}\NormalTok{, }\DecValTok{91}\NormalTok{, }\DecValTok{93}\NormalTok{, }\DecValTok{65}\NormalTok{),}
  \AttributeTok{attendance =} \FunctionTok{c}\NormalTok{(}\DecValTok{95}\NormalTok{, }\DecValTok{88}\NormalTok{, }\DecValTok{100}\NormalTok{, }\DecValTok{72}\NormalTok{, }\DecValTok{90}\NormalTok{, }\DecValTok{85}\NormalTok{, }\DecValTok{98}\NormalTok{, }\DecValTok{60}\NormalTok{)}
\NormalTok{)}
\NormalTok{students}
\end{Highlighting}
\end{Shaded}

\begin{verbatim}
##   id  name sex stringAsFactor section midterm final attendance
## 1  1 Alice   F          FALSE       A      78    82         95
## 2  2 Bruno   M          FALSE       B      85    79         88
## 3  3  Chen   M          FALSE       A      92    96        100
## 4  4  Dana   F          FALSE       C      67    71         72
## 5  5  Elif   F          FALSE       B      74    80         90
## 6  6 Farid   M          FALSE       C      88    91         85
## 7  7 Grace   F          FALSE       A      95    93         98
## 8  8  Hiro   M          FALSE       B      59    65         60
\end{verbatim}

\textbf{B2} Display the structure, dimensions, column names, and summary
of \texttt{students}

\begin{Shaded}
\begin{Highlighting}[]
\FunctionTok{str}\NormalTok{(students)}
\end{Highlighting}
\end{Shaded}

\begin{verbatim}
## 'data.frame':    8 obs. of  8 variables:
##  $ id            : int  1 2 3 4 5 6 7 8
##  $ name          : chr  "Alice" "Bruno" "Chen" "Dana" ...
##  $ sex           : chr  "F" "M" "M" "F" ...
##  $ stringAsFactor: logi  FALSE FALSE FALSE FALSE FALSE FALSE ...
##  $ section       : chr  "A" "B" "A" "C" ...
##  $ midterm       : num  78 85 92 67 74 88 95 59
##  $ final         : num  82 79 96 71 80 91 93 65
##  $ attendance    : num  95 88 100 72 90 85 98 60
\end{verbatim}

\begin{Shaded}
\begin{Highlighting}[]
\FunctionTok{dim}\NormalTok{(students)}
\end{Highlighting}
\end{Shaded}

\begin{verbatim}
## [1] 8 8
\end{verbatim}

\begin{Shaded}
\begin{Highlighting}[]
\FunctionTok{colnames}\NormalTok{(students)}
\end{Highlighting}
\end{Shaded}

\begin{verbatim}
## [1] "id"             "name"           "sex"            "stringAsFactor"
## [5] "section"        "midterm"        "final"          "attendance"
\end{verbatim}

\begin{Shaded}
\begin{Highlighting}[]
\FunctionTok{summary}\NormalTok{(students)}
\end{Highlighting}
\end{Shaded}

\begin{verbatim}
##        id              name          sex    stringAsFactor       section 
##  Min.   :1.00   Length   :8   Length   :8   Mode :logical   Length   :8  
##  1st Qu.:2.75   N.unique :8   N.unique :2   FALSE:8         N.unique :3  
##  Median :4.50   N.blank  :0   N.blank  :0                   N.blank  :0  
##  Mean   :4.50   Min.nchar:4   Min.nchar:1                   Min.nchar:1  
##  3rd Qu.:6.25   Max.nchar:5   Max.nchar:1                   Max.nchar:1  
##  Max.   :8.00                                                            
##     midterm          final         attendance    
##  Min.   :59.00   Min.   :65.00   Min.   : 60.00  
##  1st Qu.:72.25   1st Qu.:77.00   1st Qu.: 81.75  
##  Median :81.50   Median :81.00   Median : 89.00  
##  Mean   :79.75   Mean   :82.12   Mean   : 86.00  
##  3rd Qu.:89.00   3rd Qu.:91.50   3rd Qu.: 95.75  
##  Max.   :95.00   Max.   :96.00   Max.   :100.00
\end{verbatim}

\textbf{B3} Add a column \texttt{average} = the mean of the midterm and
final for each studen, rounded to 1 decimal.

\begin{Shaded}
\begin{Highlighting}[]
\NormalTok{students}\SpecialCharTok{$}\NormalTok{average }\OtherTok{\textless{}{-}} \FunctionTok{round}\NormalTok{(}\FunctionTok{mean}\NormalTok{(}\FunctionTok{c}\NormalTok{(students}\SpecialCharTok{$}\NormalTok{midterm, students}\SpecialCharTok{$}\NormalTok{final)), }\AttributeTok{digits =} \DecValTok{1}\NormalTok{)}
\NormalTok{students}\SpecialCharTok{$}\NormalTok{average}
\end{Highlighting}
\end{Shaded}

\begin{verbatim}
## [1] 80.9 80.9 80.9 80.9 80.9 80.9 80.9 80.9
\end{verbatim}

\begin{Shaded}
\begin{Highlighting}[]
\NormalTok{students}
\end{Highlighting}
\end{Shaded}

\begin{verbatim}
##   id  name sex stringAsFactor section midterm final attendance average
## 1  1 Alice   F          FALSE       A      78    82         95    80.9
## 2  2 Bruno   M          FALSE       B      85    79         88    80.9
## 3  3  Chen   M          FALSE       A      92    96        100    80.9
## 4  4  Dana   F          FALSE       C      67    71         72    80.9
## 5  5  Elif   F          FALSE       B      74    80         90    80.9
## 6  6 Farid   M          FALSE       C      88    91         85    80.9
## 7  7 Grace   F          FALSE       A      95    93         98    80.9
## 8  8  Hiro   M          FALSE       B      59    65         60    80.9
\end{verbatim}

\textbf{B4} Add a logical column passed that is TRUE when average
\textgreater= 70 and attendance \textgreater= 75.

\begin{Shaded}
\begin{Highlighting}[]
\NormalTok{students}\SpecialCharTok{$}\NormalTok{passed }\OtherTok{\textless{}{-}}\NormalTok{ students}\SpecialCharTok{$}\NormalTok{average }\SpecialCharTok{\textgreater{}=} \DecValTok{70} \SpecialCharTok{\&}\NormalTok{ students}\SpecialCharTok{$}\NormalTok{attendance }\SpecialCharTok{\textgreater{}=} \DecValTok{75}
\NormalTok{students}\SpecialCharTok{$}\NormalTok{passed}
\end{Highlighting}
\end{Shaded}

\begin{verbatim}
## [1]  TRUE  TRUE  TRUE FALSE  TRUE  TRUE  TRUE FALSE
\end{verbatim}

\begin{Shaded}
\begin{Highlighting}[]
\NormalTok{students}
\end{Highlighting}
\end{Shaded}

\begin{verbatim}
##   id  name sex stringAsFactor section midterm final attendance average passed
## 1  1 Alice   F          FALSE       A      78    82         95    80.9   TRUE
## 2  2 Bruno   M          FALSE       B      85    79         88    80.9   TRUE
## 3  3  Chen   M          FALSE       A      92    96        100    80.9   TRUE
## 4  4  Dana   F          FALSE       C      67    71         72    80.9  FALSE
## 5  5  Elif   F          FALSE       B      74    80         90    80.9   TRUE
## 6  6 Farid   M          FALSE       C      88    91         85    80.9   TRUE
## 7  7 Grace   F          FALSE       A      95    93         98    80.9   TRUE
## 8  8  Hiro   M          FALSE       B      59    65         60    80.9  FALSE
\end{verbatim}

\textbf{B5} Using ifelse(), add a column standing with these values
based on average: . ``Honors'' if average \textgreater= 90 . ``Good'' if
average \textgreater= 75 (and below 90) . ``Needs Support'' otherwise

\begin{Shaded}
\begin{Highlighting}[]
\DocumentationTok{\#\#students$remarks \textless{}{-} ifelse(average \textgreater{}= 90, "Honors", "Good")}
\DocumentationTok{\#\#ifelse(average \textgreater{}= 75, "Good", "Needs Support")}
\end{Highlighting}
\end{Shaded}

\textbf{B6} Write a function classify(score) that uses an if / else if /
else chain and returns ``Distinction'' for 85 or above, ``Merit'' for 70
to 84, ``Pass'' for 50 to 69, and ``Fail'' below 50. Test it with
sapply(c(90, 73, 55, 40), classify) and print the result.

\begin{Shaded}
\begin{Highlighting}[]
\DocumentationTok{\#\#classify(score) \textless{}{-} if()}
\end{Highlighting}
\end{Shaded}

\subsubsection{\texorpdfstring{Part c: \texttt{rbind()} and
\texttt{cbind()}}{Part c: rbind() and cbind()}}\label{part-c-rbind-and-cbind}

\textbf{C1} Create a data frame new\_students with two students. Its
columns must match the original, including average, passed and standing,
computed the same way: * id 9, Ines, F, section C, midterm 81, final 84,
attendance 92 * id 10, Jamal, M, section A, midterm 70, final 68,
attendance 80 Use \texttt{rbind()} to append them to students, saving
the result as students\_all.

\begin{Shaded}
\begin{Highlighting}[]
\NormalTok{students\_all }\OtherTok{\textless{}{-}} \FunctionTok{rbind}\NormalTok{(students}\SpecialCharTok{$}\NormalTok{id, students}\SpecialCharTok{$}\NormalTok{section, students}\SpecialCharTok{$}\NormalTok{midterm, students}\SpecialCharTok{$}\NormalTok{final, students}\SpecialCharTok{$}\NormalTok{attendance)}
\NormalTok{students\_all}
\end{Highlighting}
\end{Shaded}

\begin{verbatim}
##      [,1] [,2] [,3]  [,4] [,5] [,6] [,7] [,8]
## [1,] "1"  "2"  "3"   "4"  "5"  "6"  "7"  "8" 
## [2,] "A"  "B"  "A"   "C"  "B"  "C"  "A"  "B" 
## [3,] "78" "85" "92"  "67" "74" "88" "95" "59"
## [4,] "82" "79" "96"  "71" "80" "91" "93" "65"
## [5,] "95" "88" "100" "72" "90" "85" "98" "60"
\end{verbatim}

\textbf{C2} Create the numeric vector quiz \textless- c(8, 7, 10, 6, 7,
9, 10, 5, 8, 6). Use cbind() to attach it to students\_all as a column
named quiz. Confirm the result is still a data frame with class().

\begin{Shaded}
\begin{Highlighting}[]
\NormalTok{quiz }\OtherTok{\textless{}{-}} \FunctionTok{c}\NormalTok{(}\DecValTok{8}\NormalTok{, }\DecValTok{7}\NormalTok{, }\DecValTok{10}\NormalTok{, }\DecValTok{6}\NormalTok{, }\DecValTok{7}\NormalTok{, }\DecValTok{9}\NormalTok{, }\DecValTok{10}\NormalTok{, }\DecValTok{5}\NormalTok{, }\DecValTok{8}\NormalTok{, }\DecValTok{6}\NormalTok{)}
\NormalTok{students\_all }\OtherTok{\textless{}{-}} \FunctionTok{cbind}\NormalTok{(students\_all, quiz)}
\end{Highlighting}
\end{Shaded}

\begin{verbatim}
## Warning in cbind(students_all, quiz): number of rows of result is not a
## multiple of vector length (arg 2)
\end{verbatim}

\begin{Shaded}
\begin{Highlighting}[]
\NormalTok{students\_all}
\end{Highlighting}
\end{Shaded}

\begin{verbatim}
##                                               quiz
## [1,] "1"  "2"  "3"   "4"  "5"  "6"  "7"  "8"  "8" 
## [2,] "A"  "B"  "A"   "C"  "B"  "C"  "A"  "B"  "7" 
## [3,] "78" "85" "92"  "67" "74" "88" "95" "59" "10"
## [4,] "82" "79" "96"  "71" "80" "91" "93" "65" "6" 
## [5,] "95" "88" "100" "72" "90" "85" "98" "60" "7"
\end{verbatim}

\begin{Shaded}
\begin{Highlighting}[]
\FunctionTok{class}\NormalTok{(students\_all)}
\end{Highlighting}
\end{Shaded}

\begin{verbatim}
## [1] "matrix" "array"
\end{verbatim}

\textbf{C3} Create the matrix
\texttt{m\ \textless{}-\ cbind(midterm\ =\ students\_all\$midterm,\ final\ =\ students\_all\$final})
and use rbind()` to append a final row containing the column means,
labelled ``mean''.

\begin{Shaded}
\begin{Highlighting}[]
\DocumentationTok{\#\#m \textless{}{-} cbind(midterm = students\_all$midterm, final = students\_all$final)}
\DocumentationTok{\#\#mean \textless{}{-} rbind(students\_all, m)}
\end{Highlighting}
\end{Shaded}

\textbf{C4} In a comment (2-3 sentences), explain what error rbind()
produces if the data frames have different column names, and show the
output if you safely catch it with tryCatch()
\texttt{tryCatch(rbind(students,\ bad),\ error\ =\ function(e)\ \{\ message("rbind\ failed:\ ",\ conditionMessage(e));\ NULL\ \})}.

\subsubsection{Part D: Subsetting}\label{part-d-subsetting}

Use students\_all (with quiz) for all questions, Print every result.

D1 Extract the name column in three ways: \$, {[}{[}``name''{]}{]}, and
{[} , ``nane''{]}.

\begin{Shaded}
\begin{Highlighting}[]
\NormalTok{students\_all }\OtherTok{\textless{}{-}}\NormalTok{ students}
\NormalTok{students\_all}\SpecialCharTok{$}\NormalTok{name}
\end{Highlighting}
\end{Shaded}

\begin{verbatim}
## [1] "Alice" "Bruno" "Chen"  "Dana"  "Elif"  "Farid" "Grace" "Hiro"
\end{verbatim}

\begin{Shaded}
\begin{Highlighting}[]
\NormalTok{students\_all[[}\StringTok{"name"}\NormalTok{]]}
\end{Highlighting}
\end{Shaded}

\begin{verbatim}
## [1] "Alice" "Bruno" "Chen"  "Dana"  "Elif"  "Farid" "Grace" "Hiro"
\end{verbatim}

\begin{Shaded}
\begin{Highlighting}[]
\NormalTok{students\_all[ , }\StringTok{"name"}\NormalTok{]}
\end{Highlighting}
\end{Shaded}

\begin{verbatim}
## [1] "Alice" "Bruno" "Chen"  "Dana"  "Elif"  "Farid" "Grace" "Hiro"
\end{verbatim}

D2 Extract rows 2 to 5 and only the columns name, midterm, final.

\begin{Shaded}
\begin{Highlighting}[]
\NormalTok{students\_all[}\DecValTok{2}\SpecialCharTok{:}\DecValTok{5}\NormalTok{, ]}
\end{Highlighting}
\end{Shaded}

\begin{verbatim}
##   id  name sex stringAsFactor section midterm final attendance average passed
## 2  2 Bruno   M          FALSE       B      85    79         88    80.9   TRUE
## 3  3  Chen   M          FALSE       A      92    96        100    80.9   TRUE
## 4  4  Dana   F          FALSE       C      67    71         72    80.9  FALSE
## 5  5  Elif   F          FALSE       B      74    80         90    80.9   TRUE
\end{verbatim}

D3 Using logical subsetting, select students withattendance' above 90.

D4 Using subset(), select students of sex ``F'' in section ``A'' or
``C'' whose final is at least 75; return only name, section and final.

D5 Use which.max() and which.min() to report the name of the student
with the highest and lowest average

D6 Sort students\_all by average in descending order, then print the top
3 rows.

\subsubsection{Part E: Basic Factors}\label{part-e-basic-factors}

El Convert section and sex in students\_all to factors with factor().
Print levels() of each.

E2 Use table() to count students per sex and per section.

Median :4.5

*Try the two-way table of sex by section. table(students\_all\$sex,
students\_allSsection). What is the output?

E3 Print as. integer() of the sex factor. In a comment, explain what the
integers represent. E4 Create a new factor sex\_label from sex using
factor() with levels = c(``F'', ``M'') and labels = ``,''Nale''). Add it
to students\_all and print its levels().

E5 Print sunmary() of sex\_label. Then, using logical subsetting and
nean(), print the mean average for each level of sex\_label.

\end{document}
